%% bare_jrnl.tex
%% V1.4b
%% 2015/08/26
%% by Michael Shell
%% see http://www.michaelshell.org/
%% for current contact information.
%%
%% This is a skeleton file demonstrating the use of IEEEtran.cls
%% (requires IEEEtran.cls version 1.8b or later) with an IEEE
%% journal paper.
%%
%% Support sites:
%% http://www.michaelshell.org/tex/ieeetran/
%% http://www.ctan.org/pkg/ieeetran
%% and
%% http://www.ieee.org/

%%*************************************************************************
%% Legal Notice:
%% This code is offered as-is without any warranty either expressed or
%% implied; without even the implied warranty of MERCHANTABILITY or
%% FITNESS FOR A PARTICULAR PURPOSE! 
%% User assumes all risk.
%% In no event shall the IEEE or any contributor to this code be liable for
%% any damages or losses, including, but not limited to, incidental,
%% consequential, or any other damages, resulting from the use or misuse
%% of any information contained here.
%%
%% All comments are the opinions of their respective authors and are not
%% necessarily endorsed by the IEEE.
%%
%% This work is distributed under the LaTeX Project Public License (LPPL)
%% ( http://www.latex-project.org/ ) version 1.3, and may be freely used,
%% distributed and modified. A copy of the LPPL, version 1.3, is included
%% in the base LaTeX documentation of all distributions of LaTeX released
%% 2003/12/01 or later.
%% Retain all contribution notices and credits.
%% ** Modified files should be clearly indicated as such, including  **
%% ** renaming them and changing author support contact information. **
%%*************************************************************************


% *** Authors should verify (and, if needed, correct) their LaTeX system  ***
% *** with the testflow diagnostic prior to trusting their LaTeX platform ***
% *** with production work. The IEEE's font choices and paper sizes can   ***
% *** trigger bugs that do not appear when using other class files.       ***                          ***
% The testflow support page is at:
% http://www.michaelshell.org/tex/testflow/



\documentclass[journal]{IEEEtran}
%
% If IEEEtran.cls has not been installed into the LaTeX system files,
% manually specify the path to it like:
% \documentclass[journal]{../sty/IEEEtran}





% Some very useful LaTeX packages include:
% (uncomment the ones you want to load)


% *** MISC UTILITY PACKAGES ***
%
%\usepackage{ifpdf}
\usepackage{url}
% Heiko Oberdiek's ifpdf.sty is very useful if you need conditional
% compilation based on whether the output is pdf or dvi.
% usage:
% \ifpdf
%   % pdf code
% \else
%   % dvi code
% \fi
% The latest version of ifpdf.sty can be obtained from:
% http://www.ctan.org/pkg/ifpdf
% Also, note that IEEEtran.cls V1.7 and later provides a builtin
% \ifCLASSINFOpdf conditional that works the same way.
% When switching from latex to pdflatex and vice-versa, the compiler may
% have to be run twice to clear warning/error messages.






% *** CITATION PACKAGES ***
%
%\usepackage{cite}
% cite.sty was written by Donald Arseneau
% V1.6 and later of IEEEtran pre-defines the format of the cite.sty package
% \cite{} output to follow that of the IEEE. Loading the cite package will
% result in citation numbers being automatically sorted and properly
% "compressed/ranged". e.g., [1], [9], [2], [7], [5], [6] without using
% cite.sty will become [1], [2], [5]--[7], [9] using cite.sty. cite.sty's
% \cite will automatically add leading space, if needed. Use cite.sty's
% noadjust option (cite.sty V3.8 and later) if you want to turn this off
% such as if a citation ever needs to be enclosed in parenthesis.
% cite.sty is already installed on most LaTeX systems. Be sure and use
% version 5.0 (2009-03-20) and later if using hyperref.sty.
% The latest version can be obtained at:
% http://www.ctan.org/pkg/cite
% The documentation is contained in the cite.sty file itself.






% *** GRAPHICS RELATED PACKAGES ***
%
\usepackage{svg}      % para \includesvg
\ifCLASSINFOpdf
  % \usepackage[pdftex]{graphicx}
  % declare the path(s) where your graphic files are
  % \graphicspath{{../pdf/}{../jpeg/}}
  % and their extensions so you won't have to specify these with
  % every instance of \includegraphics
  % \DeclareGraphicsExtensions{.pdf,.jpeg,.png}
\else
  % or other class option (dvipsone, dvipdf, if not using dvips). graphicx
  % will default to the driver specified in the system graphics.cfg if no
  % driver is specified.
  % \usepackage[dvips]{graphicx}
  % declare the path(s) where your graphic files are
  % \graphicspath{{../eps/}}
  % and their extensions so you won't have to specify these with
  % every instance of \includegraphics
  % \DeclareGraphicsExtensions{.eps}
\fi
% graphicx was written by David Carlisle and Sebastian Rahtz. It is
% required if you want graphics, photos, etc. graphicx.sty is already
% installed on most LaTeX systems. The latest version and documentation
% can be obtained at: 
% http://www.ctan.org/pkg/graphicx
% Another good source of documentation is "Using Imported Graphics in
% LaTeX2e" by Keith Reckdahl which can be found at:
% http://www.ctan.org/pkg/epslatex
%
% latex, and pdflatex in dvi mode, support graphics in encapsulated
% postscript (.eps) format. pdflatex in pdf mode supports graphics
% in .pdf, .jpeg, .png and .mps (metapost) formats. Users should ensure
% that all non-photo figures use a vector format (.eps, .pdf, .mps) and
% not a bitmapped formats (.jpeg, .png). The IEEE frowns on bitmapped formats
% which can result in "jaggedy"/blurry rendering of lines and letters as
% well as large increases in file sizes.
%
% You can find documentation about the pdfTeX application at:
% http://www.tug.org/applications/pdftex





% *** MATH PACKAGES ***
%
%\usepackage{amsmath}
% A popular package from the American Mathematical Society that provides
% many useful and powerful commands for dealing with mathematics.
%
% Note that the amsmath package sets \interdisplaylinepenalty to 10000
% thus preventing page breaks from occurring within multiline equations. Use:
%\interdisplaylinepenalty=2500
% after loading amsmath to restore such page breaks as IEEEtran.cls normally
% does. amsmath.sty is already installed on most LaTeX systems. The latest
% version and documentation can be obtained at:
% http://www.ctan.org/pkg/amsmath





% *** SPECIALIZED LIST PACKAGES ***
%
%\usepackage{algorithmic}
% algorithmic.sty was written by Peter Williams and Rogerio Brito.
% This package provides an algorithmic environment fo describing algorithms.
% You can use the algorithmic environment in-text or within a figure
% environment to provide for a floating algorithm. Do NOT use the algorithm
% floating environment provided by algorithm.sty (by the same authors) or
% algorithm2e.sty (by Christophe Fiorio) as the IEEE does not use dedicated
% algorithm float types and packages that provide these will not provide
% correct IEEE style captions. The latest version and documentation of
% algorithmic.sty can be obtained at:
% http://www.ctan.org/pkg/algorithms
% Also of interest may be the (relatively newer and more customizable)
% algorithmicx.sty package by Szasz Janos:
% http://www.ctan.org/pkg/algorithmicx




% *** ALIGNMENT PACKAGES ***
%
%\usepackage{array}
% Frank Mittelbach's and David Carlisle's array.sty patches and improves
% the standard LaTeX2e array and tabular environments to provide better
% appearance and additional user controls. As the default LaTeX2e table
% generation code is lacking to the point of almost being broken with
% respect to the quality of the end results, all users are strongly
% advised to use an enhanced (at the very least that provided by array.sty)
% set of table tools. array.sty is already installed on most systems. The
% latest version and documentation can be obtained at:
% http://www.ctan.org/pkg/array


% IEEEtran contains the IEEEeqnarray family of commands that can be used to
% generate multiline equations as well as matrices, tables, etc., of high
% quality.




% *** SUBFIGURE PACKAGES ***
%\ifCLASSOPTIONcompsoc
%  \usepackage[caption=false,font=normalsize,labelfont=sf,textfont=sf]{subfig}
%\else
%  \usepackage[caption=false,font=footnotesize]{subfig}
%\fi
% subfig.sty, written by Steven Douglas Cochran, is the modern replacement
% for subfigure.sty, the latter of which is no longer maintained and is
% incompatible with some LaTeX packages including fixltx2e. However,
% subfig.sty requires and automatically loads Axel Sommerfeldt's caption.sty
% which will override IEEEtran.cls' handling of captions and this will result
% in non-IEEE style figure/table captions. To prevent this problem, be sure
% and invoke subfig.sty's "caption=false" package option (available since
% subfig.sty version 1.3, 2005/06/28) as this is will preserve IEEEtran.cls
% handling of captions.
% Note that the Computer Society format requires a larger sans serif font
% than the serif footnote size font used in traditional IEEE formatting
% and thus the need to invoke different subfig.sty package options depending
% on whether compsoc mode has been enabled.
%
% The latest version and documentation of subfig.sty can be obtained at:
% http://www.ctan.org/pkg/subfig




% *** FLOAT PACKAGES ***
%
%\usepackage{fixltx2e}
% fixltx2e, the successor to the earlier fix2col.sty, was written by
% Frank Mittelbach and David Carlisle. This package corrects a few problems
% in the LaTeX2e kernel, the most notable of which is that in current
% LaTeX2e releases, the ordering of single and double column floats is not
% guaranteed to be preserved. Thus, an unpatched LaTeX2e can allow a
% single column figure to be placed prior to an earlier double column
% figure.
% Be aware that LaTeX2e kernels dated 2015 and later have fixltx2e.sty's
% corrections already built into the system in which case a warning will
% be issued if an attempt is made to load fixltx2e.sty as it is no longer
% needed.
% The latest version and documentation can be found at:
% http://www.ctan.org/pkg/fixltx2e


%\usepackage{stfloats}
% stfloats.sty was written by Sigitas Tolusis. This package gives LaTeX2e
% the ability to do double column floats at the bottom of the page as well
% as the top. (e.g., "\begin{figure*}[!b]" is not normally possible in
% LaTeX2e). It also provides a command:
%\fnbelowfloat
% to enable the placement of footnotes below bottom floats (the standard
% LaTeX2e kernel puts them above bottom floats). This is an invasive package
% which rewrites many portions of the LaTeX2e float routines. It may not work
% with other packages that modify the LaTeX2e float routines. The latest
% version and documentation can be obtained at:
% http://www.ctan.org/pkg/stfloats
% Do not use the stfloats baselinefloat ability as the IEEE does not allow
% \baselineskip to stretch. Authors submitting work to the IEEE should note
% that the IEEE rarely uses double column equations and that authors should try
% to avoid such use. Do not be tempted to use the cuted.sty or midfloat.sty
% packages (also by Sigitas Tolusis) as the IEEE does not format its papers in
% such ways.
% Do not attempt to use stfloats with fixltx2e as they are incompatible.
% Instead, use Morten Hogholm'a dblfloatfix which combines the features
% of both fixltx2e and stfloats:
%
% \usepackage{dblfloatfix}
% The latest version can be found at:
% http://www.ctan.org/pkg/dblfloatfix




%\ifCLASSOPTIONcaptionsoff
%  \usepackage[nomarkers]{endfloat}
% \let\MYoriglatexcaption\caption
% \renewcommand{\caption}[2][\relax]{\MYoriglatexcaption[#2]{#2}}
%\fi
% endfloat.sty was written by James Darrell McCauley, Jeff Goldberg and 
% Axel Sommerfeldt. This package may be useful when used in conjunction with 
% IEEEtran.cls'  captionsoff option. Some IEEE journals/societies require that
% submissions have lists of figures/tables at the end of the paper and that
% figures/tables without any captions are placed on a page by themselves at
% the end of the document. If needed, the draftcls IEEEtran class option or
% \CLASSINPUTbaselinestretch interface can be used to increase the line
% spacing as well. Be sure and use the nomarkers option of endfloat to
% prevent endfloat from "marking" where the figures would have been placed
% in the text. The two hack lines of code above are a slight modification of
% that suggested by in the endfloat docs (section 8.4.1) to ensure that
% the full captions always appear in the list of figures/tables - even if
% the user used the short optional argument of \caption[]{}.
% IEEE papers do not typically make use of \caption[]'s optional argument,
% so this should not be an issue. A similar trick can be used to disable
% captions of packages such as subfig.sty that lack options to turn off
% the subcaptions:
% For subfig.sty:
% \let\MYorigsubfloat\subfloat
% \renewcommand{\subfloat}[2][\relax]{\MYorigsubfloat[]{#2}}
% However, the above trick will not work if both optional arguments of
% the \subfloat command are used. Furthermore, there needs to be a
% description of each subfigure *somewhere* and endfloat does not add
% subfigure captions to its list of figures. Thus, the best approach is to
% avoid the use of subfigure captions (many IEEE journals avoid them anyway)
% and instead reference/explain all the subfigures within the main caption.
% The latest version of endfloat.sty and its documentation can obtained at:
% http://www.ctan.org/pkg/endfloat
%
% The IEEEtran \ifCLASSOPTIONcaptionsoff conditional can also be used
% later in the document, say, to conditionally put the References on a 
% page by themselves.




% *** PDF, URL AND HYPERLINK PACKAGES ***
%
%\usepackage{url}
% url.sty was written by Donald Arseneau. It provides better support for
% handling and breaking URLs. url.sty is already installed on most LaTeX
% systems. The latest version and documentation can be obtained at:
% http://www.ctan.org/pkg/url
% Basically, \url{my_url_here}.




% *** Do not adjust lengths that control margins, column widths, etc. ***
% *** Do not use packages that alter fonts (such as pslatex).         ***
% There should be no need to do such things with IEEEtran.cls V1.6 and later.
% (Unless specifically asked to do so by the journal or conference you plan
% to submit to, of course. )


% correct bad hyphenation here
\hyphenation{op-tical net-works semi-conduc-tor}


\begin{document}
%
% paper title
% Titles are generally capitalized except for words such as a, an, and, as,
% at, but, by, for, in, nor, of, on, or, the, to and up, which are usually
% not capitalized unless they are the first or last word of the title.
% Linebreaks \\ can be used within to get better formatting as desired.
% Do not put math or special symbols in the title.
\title{Bare Demo of IEEEtran.cls\\ for IEEE Journals}
%
%
% author names and IEEE memberships
% note positions of commas and nonbreaking spaces ( ~ ) LaTeX will not break
% a structure at a ~ so this keeps an author's name from being broken across
% two lines.
% use \thanks{} to gain access to the first footnote area
% a separate \thanks must be used for each paragraph as LaTeX2e's \thanks
% was not built to handle multiple paragraphs
%

\author{Michael~Shell,~\IEEEmembership{Member,~IEEE,}
        John~Doe,~\IEEEmembership{Fellow,~OSA,}
        and~Jane~Doe,~\IEEEmembership{Life~Fellow,~IEEE}% <-this % stops a space
\thanks{M. Shell was with the Department
of Electrical and Computer Engineering, Georgia Institute of Technology, Atlanta,
GA, 30332 USA e-mail: (see http://www.michaelshell.org/contact.html).}% <-this % stops a space
\thanks{J. Doe and J. Doe are with Anonymous University.}% <-this % stops a space
\thanks{Manuscript received April 19, 2005; revised August 26, 2015.}}

% note the % following the last \IEEEmembership and also \thanks - 
% these prevent an unwanted space from occurring between the last author name
% and the end of the author line. i.e., if you had this:
% 
% \author{....lastname \thanks{...} \thanks{...} }
%                     ^------------^------------^----Do not want these spaces!
%
% a space would be appended to the last name and could cause every name on that
% line to be shifted left slightly. This is one of those "LaTeX things". For
% instance, "\textbf{A} \textbf{B}" will typeset as "A B" not "AB". To get
% "AB" then you have to do: "\textbf{A}\textbf{B}"
% \thanks is no different in this regard, so shield the last } of each \thanks
% that ends a line with a % and do not let a space in before the next \thanks.
% Spaces after \IEEEmembership other than the last one are OK (and needed) as
% you are supposed to have spaces between the names. For what it is worth,
% this is a minor point as most people would not even notice if the said evil
% space somehow managed to creep in.



% The paper headers
\markboth{Journal of \LaTeX\ Class Files,~Vol.~14, No.~8, August~2015}%
{Shell \MakeLowercase{\textit{et al.}}: Bare Demo of IEEEtran.cls for IEEE Journals}
% The only time the second header will appear is for the odd numbered pages
% after the title page when using the twoside option.
% 
% *** Note that you probably will NOT want to include the author's ***
% *** name in the headers of peer review papers.                   ***
% You can use \ifCLASSOPTIONpeerreview for conditional compilation here if
% you desire.




% If you want to put a publisher's ID mark on the page you can do it like
% this:
%\IEEEpubid{0000--0000/00\$00.00~\copyright~2015 IEEE}
% Remember, if you use this you must call \IEEEpubidadjcol in the second
% column for its text to clear the IEEEpubid mark.



% use for special paper notices
%\IEEEspecialpapernotice{(Invited Paper)}




% make the title area
\maketitle

% As a general rule, do not put math, special symbols or citations
% in the abstract or keywords.
\begin{abstract}
The abstract goes here.
\end{abstract}

% Note that keywords are not normally used for peerreview papers.
\begin{IEEEkeywords}
IEEE, IEEEtran, journal, \LaTeX, paper, template.
\end{IEEEkeywords}






% For peer review papers, you can put extra information on the cover
% page as needed:
% \ifCLASSOPTIONpeerreview
% \begin{center} \bfseries EDICS Category: 3-BBND \end{center}
% \fi
%
% For peerreview papers, this IEEEtran command inserts a page break and
% creates the second title. It will be ignored for other modes.
\IEEEpeerreviewmaketitle



\section{Marco teórico}
% The very first letter is a 2 line initial drop letter followed
% by the rest of the first word in caps.
% 
% form to use if the first word consists of a single letter:
% \IEEEPARstart{A}{demo} file is ....
% 
% form to use if you need the single drop letter followed by
% normal text (unknown if ever used by the IEEE):
% \IEEEPARstart{A}{}demo file is ....
% 
% Some journals put the first two words in caps:
% \IEEEPARstart{T}{his demo} file is ....
\subsection{K8s Networking ¿cómo se comunican las cosas dentro y fuera del clúster?}

\begin{figure}[htbp]
  \centering
  \includesvg[width=\linewidth]{./fig/kubernetes-cluster-network}
  \caption{Topología de red de un clúster de Kubernetes (adaptado de \cite{ref2}).}
  \label{fig:k8s-cluster-network}
\end{figure}

\medskip
\subsubsection{Mapa mental de redes en Kubernetes (analogía de ciudad)}

Un clúster puede conceptualizarse como una ciudad organizada: los \emph{Pods} son las “viviendas” donde se ejecuta la aplicación; los \emph{Nodes} agrupan múltiples viviendas y proveen recursos comunes; los \emph{Services} operan como una central con un número estable para localizar réplicas; y \emph{Ingress/Gateway} implementan el acceso controlado desde el exterior (Internet), actuando como portería y punto de gobierno del tráfico.

\medskip
\subsubsection{Pods = casas (identidad y direcciones)}

Kubernetes adopta el modelo \emph{IP-por-Pod}: cada Pod obtiene su propia dirección IP, habilitando comunicación directa entre cargas sin traductores intermedios. Dado que la identidad del Pod es efímera —las réplicas se crean, destruyen y sustituyen de manera dinámica— la plataforma requiere abstracciones estables para descubrimiento y direccionamiento consistentes hacia los destinos correctos [2].

\medskip
\subsubsection{Nodes = barrios (agrupación y recursos comunes)}

Un Node hospeda múltiples Pods y suministra CPU, memoria y conectividad. En este nivel se sostienen las rutas y los mecanismos de reenvío que aseguran la entrega de paquetes al Pod adecuado, conforme al modelo de red del clúster (alcance entre Pods, no superposición de rangos, conectividad \emph{east--west}) descrito por Kubernetes [2].

\medskip
\subsubsection{Services = central telefónica (número estable y reparto)}

Un Service expone un identificador estable (p.,ej., \emph{ClusterIP} o IP virtual) que funciona como “número de central” para un conjunto de Pods equivalentes. El cliente se conecta a ese número y la plataforma balancea hacia una de las réplicas disponibles, desacoplando a los consumidores de los cambios en el ciclo de vida de los Pods [3], [13].
Para mantener actualizada la relación de destinos, Kubernetes emplea \emph{EndpointSlices}, que enumeran direcciones y puertos de los Pods que respaldan un Service, mejorando la escalabilidad y reduciendo la sobrecarga de control frente a los \emph{Endpoints} tradicionales [1].

\medskip
\subsubsection{Ingress / Gateway = portería del conjunto (entrada ordenada desde la calle)}

El tráfico externo no accede directamente a los Pods: ingresa mediante una capa de entrada. Con \emph{Gateway API} se formalizan roles (infraestructura, operador, equipo de aplicación) y reglas L7 (host, rutas, TLS) para decidir por qué puerta se entra y a qué Service debe encaminarse cada solicitud [17]. Esta capa organiza el ingreso HTTP/HTTPS y refuerza la separación de responsabilidades sin exponer puertos arbitrarios.

\medskip
\subsubsection{Calles y cableado: CNI (cómo se conectan barrios y casas)}

La malla vial la provee \emph{Container Network Interface} (CNI): un conjunto de plugins intercambiables que conectan Pods y Nodes, asignan IPs y programan rutas. Kubernetes requiere un plugin compatible para la conectividad básica; la especificación CNI define cómo el \emph{runtime} solicita “conectar una vivienda” (interfaces, IP y rutas) y cómo liberar recursos al eliminarla [4], [18].
A modo ilustrativo, \emph{host-local} (asignación local de direcciones) y \emph{port-map} (traducción de puertos con el host) muestran que CNI cubre identidad y conectividad de red, no funciones avanzadas de capa 7 [20], [21].

\medskip
\subsubsection{La autopista previa a la portería: balanceo de carga externo}

En numerosos despliegues, el tráfico desde Internet hacia Ingress/Gateway atraviesa primero un balanceador de carga externo L3/L4. Diseños de alto rendimiento como \emph{Maglev} evidencian cómo el \emph{hash} consistente y el seguimiento de flujos contribuyen a una distribución estable y resiliente antes de alcanzar el clúster [5]. No implica que Kubernetes “use Maglev por defecto”, sino que ilustra el papel de un distribuidor previo a la portería.

\subsubsection{Namespaces = mini-barrios (señalización y orden, no muros)}

Un \emph{Namespace} es una partición lógica del clúster: aporta \emph{señalización} y \emph{orden} para nombrar y agrupar recursos, pero no constituye una barrera física. Facilita organizar por equipos o entornos (\texttt{dev}/\texttt{test}/\texttt{prod}) y evitar colisiones de nombres sin multiplicar clústeres \cite{ref2}.

\medskip
La señalización implica unicidad de nombres \emph{dentro} del namespace. Así, pueden coexistir dos servicios llamados \texttt{api} si residen en mini-barrios distintos (p.,ej., \texttt{api} en \texttt{dev} y \texttt{api} en \texttt{prod}) \cite{ref2}.

\medskip
El nombre DNS de un servicio incluye el namespace: \texttt{<servicio>.<namespace>.svc.cluster.local}. Dentro del mismo mini-barrio se admite el nombre corto (\texttt{api}); para acceder a otro, es preferible el FQDN (\texttt{api.prod.svc.cluster.local}) para evitar ambigüedades y asegurar enrutamiento predecible \cite{ref13,ref2}.

\medskip
Por defecto, los namespaces no aíslan el tráfico: la conectividad \emph{east--west} permanece abierta. Para definir \emph{quién puede hablar con quién}, se emplean \emph{NetworkPolicies} (si el CNI las soporta), especificadas con \emph{labels} y puertos. Este enfoque se alinea con \emph{Zero Trust}: permitir sólo lo necesario y aplicar verificación explícita \cite{ref4,ref24,ref25,ref26}.

\medskip
\subsubsection{Recorrido completo (resumen operativo)}

\begin{enumerate}
\item \textbf{Cliente externo} $\rightarrow$ (opcional) \textbf{balanceador} $\rightarrow$ \textbf{portería (Ingress/Gateway)} \cite{ref17}.
\item La portería aplica reglas L7 y dirige la solicitud a un \textbf{Service} (central) \cite{ref3,ref13}.
\item El Service consulta sus \textbf{EndpointSlices} (inventario de réplicas activas) y selecciona un Pod \cite{ref1}.
\item El tráfico circula por las \textbf{calles CNI} hasta la \textbf{casa (Pod)} en el \textbf{barrio (Node)} \cite{ref2,ref4,ref18}.
\end{enumerate}

En síntesis, la analogía de ciudad permite articular el \emph{K8s Networking}: Pods como viviendas (identidad IP efímera), Nodes como barrios (recursos compartidos), Services como central estable (balanceo con \emph{EndpointSlices}), Namespaces como mini-barrios para organización, e Ingress/Gateway como portería (gobierno de entrada). Todo opera sobre la malla vial provista por CNI, que entrega la conectividad fundamental del clúster [1]–[4], [13], [17], [18], [20], [21].

\subsection{CNI}

\subsubsection{¿Qué es CNI y qué resuelve?}

El \emph{Container Network Interface} (CNI) es la especificación y el conjunto de \emph{plugins} que habilitan la conectividad de red de los Pods en Kubernetes. En términos prácticos, CNI define cómo el \emph{runtime} de contenedores solicita a un plugin que “conecte” un Pod a la red: creación y unión de interfaces, asignación de direcciones IP mediante IPAM, configuración de rutas y, cuando corresponde, reglas de traducción/puertos en el \emph{host} \cite{ref18,ref4,ref21,ref20}. El ciclo de vida se basa en operaciones simples (\texttt{ADD}/\texttt{DEL}) encadenables, lo que permite combinar responsabilidades (p.\,ej., un plugin para IPAM y otro para mapeo de puertos) \cite{ref18,ref21,ref20}.

\paragraph{Funciones núcleo.}
\begin{itemize}
  \item \textbf{Identidad de red del Pod.} Otorga al Pod su propia IP (modelo \emph{IP-por-Pod}) y asegura la reachabilidad \emph{east--west} en el clúster \cite{ref2}.
  \item \textbf{Conexión y enrutamiento.} Añade interfaces virtuales, instala rutas y coordina con el \emph{host} el paso del tráfico hacia/desde cada Pod \cite{ref18,ref4}.
  \item \textbf{Gestión de direcciones (IPAM).} Reserva/entrega direcciones desde rangos predefinidos (ejemplo: \emph{host-local}) y evita colisiones \cite{ref20,ref18}.
\end{itemize}

\begin{figure}[h!]
    \centering
    \includegraphics[width=\linewidth]{./fig/cni_workflow.png}
    \caption{Workflow del Container Network Interface (CNI) mostrando las interacciones entre Kubelet, Pod y plugins CNI.}
    \label{fig:cni_workflow}
\end{figure}


\paragraph{¿Por qué importa?}
La elección de CNI (y su \emph{dataplane}) condiciona latencia, \emph{throughput} y consumo de CPU, con efectos directos sobre el sobredimensionamiento: estudios comparativos en \emph{edge}/nube muestran variaciones relevantes entre modos (túnel vs.\ ruteo nativo) y aceleraciones (\emph{fast paths}) \cite{ref7,ref8,ref11,ref22}. Además, la sintonía de parámetros como la MTU—clave en entornos con encapsulación—evita fragmentación y pérdidas de rendimiento \cite{ref12}. En conjunto, un CNI bien seleccionado y calibrado facilita automatizar políticas, sostener objetivos de nivel de servicio (\emph{SLOs}) y reducir recursos ociosos \cite{ref4,ref7,ref8,ref12}.

\subsubsection{Dataplane: overlay vs. ruteo directo}
En Kubernetes, el \emph{dataplane} de red que implementa un CNI puede seguir dos enfoques predominantes. El primero es el \emph{overlay}, que encapsula el tráfico entre nodos (p.\,ej., mediante VXLAN sobre UDP) para transportar paquetes de Pods sin requerir cambios en la red física. Su principal ventaja es la simplicidad operativa y el desacoplamiento del direccionamiento del clúster respecto al \emph{underlay}; a cambio, introduce cabeceras adicionales, posibles ajustes de MTU y cierta carga de CPU por encapsulación/descapsulación. El segundo enfoque es el ruteo nativo (\emph{ruteo directo}), en el que los Pods son alcanzables a través de rutas L3 entre nodos sin túneles; suele ofrecer menor latencia y mejor aprovechamiento de CPU, pero exige una integración clara con la red subyacente y un plan de direccionamiento coherente. Estudios comparativos en entornos \emph{edge} y nube muestran diferencias medibles de latencia, \emph{throughput} y uso de CPU entre CNIs y modos de operación, lo que refuerza la importancia de seleccionar el dataplane según el contexto de despliegue \cite{ref7,ref8}. Adicionalmente, la implementación concreta del fastpath (p.\,ej., conmutación basada en OVS o programas eBPF) condiciona el rendimiento y la observabilidad del plano de datos \cite{ref11,ref22}.

\subsubsection{Asignación de direcciones (IPAM)}
La identidad de red del Pod se materializa con una dirección IP propia, asignada por el mecanismo de \emph{IP Address Management} (IPAM) del CNI. En líneas generales, existen estrategias por nodo (cada nodo administra un rango local) o globales (un pool común), con implicaciones distintas en simplicidad operativa y uso de direcciones. El plugin \emph{host-local} ilustra un enfoque minimalista: asigna direcciones desde ficheros locales por nodo, garantizando unicidad sin dependencias externas. Un plan de direccionamiento bien definido (rangos, \emph{leases}, reutilización y limpieza) minimiza colisiones, facilita \emph{dual-stack} y reduce incidencias durante escalado o rotación de Pods \cite{ref18,ref20}.

\subsubsection{MTU y fragmentación en redes con túneles}
La encapsulación propia de los \emph{overlays} reduce la MTU efectiva: parte de los bytes del paquete se consumen en cabeceras adicionales. Si la MTU no se ajusta, pueden aparecer fragmentación, caídas de \emph{throughput} y aumentos de latencia. Las guías prácticas recomiendan calcular la MTU efectiva restando el \emph{overhead} del túnel y validar el ajuste mediante pruebas controladas; este procedimiento es especialmente relevante en despliegues con VXLAN o cifrado a nivel de red \cite{ref12}. Simplificando lo expuesto, un ajuste conservador de MTU, verificado empíricamente, suele prevenir pérdidas de rendimiento en clústeres con encapsulación.

\begin{figure}[t]
  \centering
  \includegraphics[width=\linewidth]{./fig/calico-mtu-overview.png}
  \caption{Ajuste de MTU en presencia de encapsulación (p.\,ej., VXLAN) para evitar fragmentación y pérdida de rendimiento \cite{ref12}.}
  \label{fig:calico_mtu}
\end{figure}


\subsubsection{Cómo circula el tráfico hacia los Services}
Aunque CNI proporciona conectividad entre Pods y nodos, el direccionamiento estable para el consumo de aplicaciones lo otorgan los \emph{Services} de Kubernetes. Un Service expone una Dirección IP Virtual - \emph{VIP} (p.\,ej., \emph{ClusterIP}) y se apoya en estructuras como \emph{EndpointSlices} para conocer qué Pods saludables pueden recibir tráfico. El CNI asegura que, una vez resuelta la VIP, el paquete alcance el Pod correspondiente a través de las rutas del clúster; el modo de exposición (ClusterIP/NodePort/LoadBalancer) define cómo entra y se distribuye el tráfico desde dentro o fuera del clúster. Comprender esta separación de responsabilidades (CNI = conectividad/IP; Service = \emph{VIP} y selección de destinos) es clave para razonar sobre fallos y cuellos de botella \cite{ref3,ref13,ref1,ref2}.

\subsubsection{Observabilidad sencilla a nivel CNI}
Para operar y optimizar una red de clúster es útil contar con visibilidad de “quién habla con quién”, volúmenes de tráfico y patrones temporales. Soluciones de \emph{flow visibility} a nivel CNI permiten recolectar y agregar registros de flujo (origen, destino, puertos, bytes/paquetes) y correlacionarlos con identidades lógicas (namespaces, \emph{labels}). Esta información facilita detectar rutas calientes, validar políticas y priorizar mejoras. En el ámbito de microservicios, la observabilidad de red complementa métricas de aplicación para diagnosticar lazos críticos o dependencias inesperadas \cite{ref23,ref6}.

\subsubsection{Rendimiento: qué medir y qué esperar}
La evaluación del CNI debería incluir latencia (promedio y colas P95/P99), \emph{throughput} sostenido y coste de CPU por paquete/conexión. También conviene fijar el modo del dataplane y el perfil de tráfico (tamaño de paquete, concurrencia, mezcla entrada/salida) para obtener comparaciones reproducibles. La bibliografía reporta patrones consistentes entre modos (overlay vs.\ ruteo directo) y entre \emph{fast-paths} (p.\,ej., OVS o eBPF), con variaciones significativas en entornos con recursos limitados o altas tasas de conexión \cite{ref7,ref8,ref9,ref11,ref22,ref27}. Estas métricas orientan decisiones de diseño y \emph{tuning} para cumplir SLOs de latencia y capacidad.

\subsubsection{Evitar sobredimensionamiento vía CNI}
Un CNI bien seleccionado y sintonizado ayuda a reducir recursos “por si acaso”. Tres prácticas son especialmente efectivas: (1) elegir el dataplane según el entorno (overlay para simplicidad o ruteo directo cuando se prioriza latencia/CPU), (2) ajustar MTU y validar que no haya fragmentación, y (3) automatizar \emph{NetworkPolicies} en clave de \emph{Zero Trust} para limitar tráfico innecesario y reducir trabajo de red en el plano de datos. Complementar estas decisiones con visibilidad de flujos permite cerrar el ciclo de mejora continua y sostener SLOs con menos CPU/memoria \cite{ref7,ref8,ref12,ref23,ref24,ref25}.


% 
% Here we have the typical use of a "T" for an initial drop letter
% and "HIS" in caps to complete the first word.
%\IEEEPARstart{T}{his} demo file is intended to serve as a ``starter file''
%for IEEE journal papers produced under \LaTeX\ using
%IEEEtran.cls version 1.8b and later.
% You must have at least 2 lines in the paragraph with the drop letter
% (should never be an issue)
%I wish you the best of success.

\hfill mds
 
\hfill August 26, 2015

\subsection{Subsection Heading Here}
Subsection text here.

% needed in second column of first page if using \IEEEpubid
%\IEEEpubidadjcol

\subsubsection{Subsubsection Heading Here}
Subsubsection text here.


% An example of a floating figure using the graphicx package.
% Note that \label must occur AFTER (or within) \caption.
% For figures, \caption should occur after the \includegraphics.
% Note that IEEEtran v1.7 and later has special internal code that
% is designed to preserve the operation of \label within \caption
% even when the captionsoff option is in effect. However, because
% of issues like this, it may be the safest practice to put all your
% \label just after \caption rather than within \caption{}.
%
% Reminder: the "draftcls" or "draftclsnofoot", not "draft", class
% option should be used if it is desired that the figures are to be
% displayed while in draft mode.
%
%\begin{figure}[!t]
%\centering
%\includegraphics[width=2.5in]{myfigure}
% where an .eps filename suffix will be assumed under latex, 
% and a .pdf suffix will be assumed for pdflatex; or what has been declared
% via \DeclareGraphicsExtensions.
%\caption{Simulation results for the network.}
%\label{fig_sim}
%\end{figure}

% Note that the IEEE typically puts floats only at the top, even when this
% results in a large percentage of a column being occupied by floats.


% An example of a double column floating figure using two subfigures.
% (The subfig.sty package must be loaded for this to work.)
% The subfigure \label commands are set within each subfloat command,
% and the \label for the overall figure must come after \caption.
% \hfil is used as a separator to get equal spacing.
% Watch out that the combined width of all the subfigures on a 
% line do not exceed the text width or a line break will occur.
%
%\begin{figure*}[!t]
%\centering
%\subfloat[Case I]{\includegraphics[width=2.5in]{box}%
%\label{fig_first_case}}
%\hfil
%\subfloat[Case II]{\includegraphics[width=2.5in]{box}%
%\label{fig_second_case}}
%\caption{Simulation results for the network.}
%\label{fig_sim}
%\end{figure*}
%
% Note that often IEEE papers with subfigures do not employ subfigure
% captions (using the optional argument to \subfloat[]), but instead will
% reference/describe all of them (a), (b), etc., within the main caption.
% Be aware that for subfig.sty to generate the (a), (b), etc., subfigure
% labels, the optional argument to \subfloat must be present. If a
% subcaption is not desired, just leave its contents blank,
% e.g., \subfloat[].


% An example of a floating table. Note that, for IEEE style tables, the
% \caption command should come BEFORE the table and, given that table
% captions serve much like titles, are usually capitalized except for words
% such as a, an, and, as, at, but, by, for, in, nor, of, on, or, the, to
% and up, which are usually not capitalized unless they are the first or
% last word of the caption. Table text will default to \footnotesize as
% the IEEE normally uses this smaller font for tables.
% The \label must come after \caption as always.
%
%\begin{table}[!t]
%% increase table row spacing, adjust to taste
%\renewcommand{\arraystretch}{1.3}
% if using array.sty, it might be a good idea to tweak the value of
% \extrarowheight as needed to properly center the text within the cells
%\caption{An Example of a Table}
%\label{table_example}
%\centering
%% Some packages, such as MDW tools, offer better commands for making tables
%% than the plain LaTeX2e tabular which is used here.
%\begin{tabular}{|c||c|}
%\hline
%One & Two\\
%\hline
%Three & Four\\
%\hline
%\end{tabular}
%\end{table}


% Note that the IEEE does not put floats in the very first column
% - or typically anywhere on the first page for that matter. Also,
% in-text middle ("here") positioning is typically not used, but it
% is allowed and encouraged for Computer Society conferences (but
% not Computer Society journals). Most IEEE journals/conferences use
% top floats exclusively. 
% Note that, LaTeX2e, unlike IEEE journals/conferences, places
% footnotes above bottom floats. This can be corrected via the
% \fnbelowfloat command of the stfloats package.




\section{Conclusion}
The conclusion goes here.





% if have a single appendix:
%\appendix[Proof of the Zonklar Equations]
% or
%\appendix  % for no appendix heading
% do not use \section anymore after \appendix, only \section*
% is possibly needed

% use appendices with more than one appendix
% then use \section to start each appendix
% you must declare a \section before using any
% \subsection or using \label (\appendices by itself
% starts a section numbered zero.)
%


\appendices
\section{Proof of the First Zonklar Equation}
Appendix one text goes here.

% you can choose not to have a title for an appendix
% if you want by leaving the argument blank
\section{}
Appendix two text goes here.


% use section* for acknowledgment
\section*{Acknowledgment}


The authors would like to thank...


% Can use something like this to put references on a page
% by themselves when using endfloat and the captionsoff option.
\ifCLASSOPTIONcaptionsoff
  \newpage
\fi



% trigger a \newpage just before the given reference
% number - used to balance the columns on the last page
% adjust value as needed - may need to be readjusted if
% the document is modified later
%\IEEEtriggeratref{8}
% The "triggered" command can be changed if desired:
%\IEEEtriggercmd{\enlargethispage{-5in}}

% references section

% can use a bibliography generated by BibTeX as a .bbl file
% BibTeX documentation can be easily obtained at:
% http://mirror.ctan.org/biblio/bibtex/contrib/doc/
% The IEEEtran BibTeX style support page is at:
% http://www.michaelshell.org/tex/ieeetran/bibtex/
%\bibliographystyle{IEEEtran}
% argument is your BibTeX string definitions and bibliography database(s)
%\bibliography{IEEEabrv,../bib/paper}
%
% <OR> manually copy in the resultant .bbl file
% set second argument of \begin to the number of references
% (used to reserve space for the reference number labels box)
\begin{thebibliography}{1}

\bibitem[1]{ref1}
``EndpointSlices,'' Kubernetes. [En línea]. Disponible en: \url{https://kubernetes.io/docs/concepts/services-networking/endpoint-slices/}. [Consultado: 21-ago-2025].

\bibitem[2]{ref2}
``Cluster networking,'' Kubernetes. [En línea]. Disponible en: \url{https://kubernetes.io/docs/concepts/cluster-administration/networking/}. [Consultado: 21-ago-2025].

\bibitem[3]{ref3}
``Virtual IPs and Service proxies,'' Kubernetes. [En línea]. Disponible en: \url{https://kubernetes.io/docs/reference/networking/virtual-ips/}. [Consultado: 21-ago-2025].

\bibitem[4]{ref4}
``Network plugins,'' Kubernetes. [En línea]. Disponible en: \url{https://kubernetes.io/docs/concepts/extend-kubernetes/compute-storage-net/network-plugins/}. [Consultado: 21-ago-2025].

\bibitem[5]{ref5}
D. E. Eisenbud \emph{et al.}, ``Maglev: A fast and reliable software network load balancer,'' \emph{USENIX NSDI}, 2016. [En línea]. Disponible en: \url{https://www.usenix.org/sites/default/files/nsdi16-paper-eisenbud.pdf}. [Consultado: 21-ago-2025].

\bibitem[6]{ref6}
F. Gomes, P. Rego, y F. Trinta, ``A systematic mapping study on observability of microservices-based applications: fundamentals, classifications, and challenges,'' \emph{Computing}, vol. 107, núm. 9, 2025.

\bibitem[7]{ref7}
Z. Kang, K. An, A. Gokhale, y P. Pazandak, ``A comprehensive performance evaluation of different Kubernetes CNI plugins for edge-based and containerized publish/subscribe applications,'' Vanderbilt University. [En línea]. Disponible en: \url{https://www.dre.vanderbilt.edu/~gokhale/WWW/papers/IC2E21_CNI_Eval.pdf}. [Consultado: 21-ago-2025].

\bibitem[8]{ref8}
G. Koukis, S. Skaperas, I. A. Kapetanidou, L. Mamatas, y V. Tsaoussidis, ``Performance evaluation of Kubernetes networking approaches across constraint edge environments,'' \emph{arXiv} [cs.NI], 2024.

\bibitem[9]{ref9}
MDPI. [En línea]. Disponible en: \url{https://www.mdpi.com/2079-9292/13/19/3972}. [Consultado: 21-ago-2025].

\bibitem[10]{ref10}
best practices, ``Elevating Kubernetes network security through Cilium deployment,'' FH Joanneum. [En línea]. Disponible en: \url{https://epub.fh-joanneum.at/obvfhjhs/download/pdf/11499653}. [Consultado: 21-ago-2025].

\bibitem[11]{ref11}
W. Tu, Y.-H. Wei, G. Antichi, y B. Pfaff, ``Revisiting the Open vSwitch dataplane ten years later,'' en \emph{Proceedings of the ACM SIGCOMM 2021 Conference}, 2021.

\bibitem[12]{ref12}
``Configure MTU to maximize network performance,'' Tigera (Calico). [En línea]. Disponible en: \url{https://docs.tigera.io/calico/latest/networking/configuring/mtu}. [Consultado: 21-ago-2025].

\bibitem[13]{ref13}
``Service,'' Kubernetes. [En línea]. Disponible en: \url{https://kubernetes.io/docs/concepts/services-networking/service/}. [Consultado: 21-ago-2025].

\bibitem[14]{ref14}
Polito.it. [En línea]. Disponible en: \url{https://webthesis.biblio.polito.it/28635/1/tesi.pdf}. [Consultado: 21-ago-2025].

\bibitem[15]{ref15}
Master’s thesis, ``Kubernetes networking: Comparative insights into API gateways and service mesh implementations,'' Aalto University. [En línea]. Disponible en: \url{https://aaltodoc.aalto.fi/bitstreams/000efe1b-5f19-494c-a8c7-3fdb1ac62e74/download}. [Consultado: 21-ago-2025].

\bibitem[16]{ref16}
M. Kotenko, D. Moskalyk, V. Kovach, y V. Osadchyi, ``Navigating the challenges and best practices in securing microservices architecture,'' \emph{CEUR-WS}. [En línea]. Disponible en: \url{https://ceur-ws.org/Vol-3826/paper1.pdf}. [Consultado: 21-ago-2025].

\bibitem[17]{ref17}
``Roles and personas — Kubernetes Gateway API,'' Kubernetes SIGs. [En línea]. Disponible en: \url{https://gateway-api.sigs.k8s.io/concepts/roles-and-personas/}. [Consultado: 21-ago-2025].

\bibitem[18]{ref18}
``CNI,'' \emph{cni.dev}. [En línea]. Disponible en: \url{https://www.cni.dev/docs/spec/}. [Consultado: 21-ago-2025].

\bibitem[19]{ref19}
IEEE, ``IEEE 802.1Q.'' [En línea]. Disponible en: \url{https://standards.ieee.org/ieee/802.1Q/6844/}. [Consultado: 21-ago-2025].

\bibitem[20]{ref20}
``host-local IP address management plugin,'' \emph{cni.dev}. [En línea]. Disponible en: \url{https://www.cni.dev/plugins/current/ipam/host-local/}. [Consultado: 21-ago-2025].

\bibitem[21]{ref21}
``Port-mapping plugin,'' \emph{cni.dev}. [En línea]. Disponible en: \url{https://www.cni.dev/plugins/current/meta/portmap/}. [Consultado: 21-ago-2025].

\bibitem[22]{ref22}
B. Pfaff \emph{et al.}, ``The design and implementation of Open vSwitch,'' \emph{USENIX NSDI}, 2015. [En línea]. Disponible en: \url{https://www.usenix.org/system/files/conference/nsdi15/nsdi15-paper-pfaff.pdf}. [Consultado: 21-ago-2025].

\bibitem[23]{ref23}
``Antrea — Network Flow Visibility,'' Antrea.io. [En línea]. Disponible en: \url{https://antrea.io/docs/v1.0.0/docs/network-flow-visibility/}. [Consultado: 21-ago-2025].

\bibitem[24]{ref24}
S. Rose, O. Borchert, S. Mitchell, y S. Connelly, ``Zero Trust Architecture,'' NIST Special Publication, 2020.

\bibitem[25]{ref25}
O. Borchert, G. Howell, A. Kerman, S. Rose, y M. Souppaya, ``Implementing a Zero Trust Architecture: High-level document,'' NIST, Gaithersburg, MD, 2025.

\bibitem[26]{ref26}
R. Chandramouli y Z. Butcher, ``Building secure microservices-based applications using service-mesh architecture,'' NIST, Gaithersburg, MD, 2020.

\bibitem[27]{ref27}
J. Zhang, P. Chen, Z. He, H. Chen, y X. Li, ``Real-time intrusion detection and prevention with neural network in kernel using eBPF,'' en \emph{Proc. 54th Annual IEEE/IFIP Int. Conf. on Dependable Systems and Networks (DSN)}, 2024, pp. 416--428.

\bibitem[28]{ref28} 
``CNI Essentials: Kubernetes Networking under the Hood,'' Tetrate. [En línea]. Disponible en: \url{https://tetrate.io/blog/kubernetes-networking}. [Consultado: 02-sep-2025].
\end{thebibliography}

% biography section
% 
% If you have an EPS/PDF photo (graphicx package needed) extra braces are
% needed around the contents of the optional argument to biography to prevent
% the LaTeX parser from getting confused when it sees the complicated
% \includegraphics command within an optional argument. (You could create
% your own custom macro containing the \includegraphics command to make things
% simpler here.)
%\begin{IEEEbiography}[{\includegraphics[width=1in,height=1.25in,clip,keepaspectratio]{mshell}}]{Michael Shell}
% or if you just want to reserve a space for a photo:

\begin{IEEEbiography}{Michael Shell}
Biography text here.
\end{IEEEbiography}

% if you will not have a photo at all:
\begin{IEEEbiographynophoto}{John Doe}
Biography text here.
\end{IEEEbiographynophoto}

% insert where needed to balance the two columns on the last page with
% biographies
%\newpage

\begin{IEEEbiographynophoto}{Jane Doe}
Biography text here.
\end{IEEEbiographynophoto}

% You can push biographies down or up by placing
% a \vfill before or after them. The appropriate
% use of \vfill depends on what kind of text is
% on the last page and whether or not the columns
% are being equalized.

%\vfill

% Can be used to pull up biographies so that the bottom of the last one
% is flush with the other column.
%\enlargethispage{-5in}



% that's all folks
\end{document}


